\documentclass[a4paper,12pt]{article}
\usepackage{docsTemplate}
\usepackage{longtable}

\setTitle{Norme di Progetto}
\setAuthors{Nenad Raduolovic}
\setVerificators{}
\setApprovation{}
\setVersion{0.3}
\setType{Interno}
\setDestination{Prof. Tullio Vardanega \\ & Prof. Riccardo Cardin}



\begin{document}

\include{docTitlepage}

\tableofcontents
\newpage

\section*{Registro delle modifiche}
\addcontentsline{toc}{section}{Registro delle modifiche}

\rowcolors{2}{white}{lightgray}
\begin{tabularx}{\textwidth}{|l|l|X|l|X|X|}
    \hline
    \textbf{Vers.} & \textbf{Data} & \textbf{Autore} & \textbf{Ruolo} & \textbf{Verifica} & \textbf{Oggetto} \\
    \hline
    \textbf{0.3} & 07/11/2025 & Nenad Radulovic & Analista & N/A & Inizio sezione documentazione, versionamento verbali \\
    \textbf{0.2} & 26/10/2025 & Nenad Radulovic & Da aggiungere & N/A & Pianificazione impegni, strumenti finora utilizzati \\
    \textbf{0.1} & 24/10/2025 & Nenad Radulovic & Da aggiungere & N/A & Impostazione documento, introduzione, inizio stesura processi organizzativi \\
    \hline
\end{tabularx}

\newpage
\section{Introduzione}
\subsection{Scopo del Documento}
Questo documento ha lo scopo di descrivere i metodi e le regole che il gruppo seguirà durante tutto lo sviluppo del progetto, al fine di ottenere un risultato efficace e di evitare problemi di organizzazione. Sono presenti nel documento anche i riferimenti a tutte le tecnologie e strumenti che verranno usati, insieme ai motivi per cui si è deciso di adottarli.

\newpage
\section{Processi Organizzativi}


\subsection{Gestione delle Risorse Umane}


\subsection{Comunicazione}
In questo capitolo vengono descritte tutte le modalità di comunicazione, sia sincrone che asincrone, che il gruppo adotterà per mantenere il contatto tra i propri membri e con i soggetti esterni.


\subsubsection{Comunicazioni Interne}
Per la comunicazione tra i membri del gruppo si è deciso di adottare le seguenti modalità sincrone:

\begin{itemize}
    \item \textbf{Incontri in presenza}: verranno organizzati, in base alle disponibilità dei membri, degli incontri in presenza con tutti membri del gruppo.

    \item \textbf{\underline{\nameref{strumenti:discord}}}: è la principale piattaforma che verrà usata per le riunioni da remoto. Consente di effettuare sia chiamate vocali che videochiamate di gruppo, permettendo di avere anche più canali vocali attivi contemporaneamente con la possibilità di cambiare velocemente il canale in cui si è presenti. Questo permette al gruppo di dividersi eventualmente in gruppi più piccoli in caso di necessità. Offre inoltre una chat per la condivisione di messaggi testuali o multimediali.
\end{itemize}

In supporto alle comunicazioni sincrone appena descritte, verranno utilizzate se seguenti piattaforme per le comunicazioni asincrone:
\begin{itemize}
    \item \textbf{\underline{\nameref{strumenti:whatsapp}}}: È stato creato un gruppo Whatsapp dedicato al progetto, nel quale sono presenti tutti i membri del gruppo. Si è deciso all'unanimità di utilizzare questo strumento rispetto ad altri vista la sua semplicità, versatilità e la maggiore familiarità che tutti i componenti del gruppo hanno con esso. Lo scopo di questo gruppo è quello di contenere le discussioni del gruppo che non richiedono di avere una tracciabilità.
    \item \textbf{\underline{\nameref{strumenti:telegram}}}: è stato creato un gruppo Telegram all'interno del quale sono stati creati una serie di canali (chat), ognuno dei quali è dedicato a un argomento specifico. Questo strumento ha lo scopo di aiutare i membri del gruppo nei loro compiti in quanto facilita la ricerca delle discussioni e del materiale condiviso in passato. Nel gruppo verranno inoltrare le decisioni prese, le attività in corso e la suddivisione corrente dei ruoli e degli incarichi (ovviamente tutte queste informazioni sono presenti anche nei verbali degli incontri, il gruppo serve a darne un accesso più centralizzato, sintetizzato e a portata di mano).
        Sono presenti i seguenti canali:
        \begin{itemize} \label{sec:canaliTelegram}
            \item \textbf{Generale}
            \item \textbf{To do in corso}: contiene allo stato attuala gli incarichi che sta svolgendo ogni membro
            \item \textbf{Registrazioni}: contiene le registrazioni audio delle riunioni interne
            \item \textbf{Riunioni}: contiene gli appunti fatti durante le riunioni
            \item \textbf{Verbali}: contiene le copie dei verbali
            \item \textbf{Link utili}: contiene link a risorse che alcuni membri hanno ritenuto utili nello svolgimento delle loro attività
            \item \textbf{Back End}: chat dedicata a discussioni riguardanti lo sviluppo Back End del progetto
            \item \textbf{Front End}: chat dedicata a discussioni riguardanti lo sviluppo Front End del progetto
            \item \textbf{Diari di Bordo}: contiene una copia delle presentazioni dei diari di bordo e eventuali domande che possono venire inserite per il prossimo diario di bordo
            \item \textbf{Dubbi}: chat in cui condividono dubbi che potrebbero impattare in modo significativo lo svolgimento del progetto. Generalmente sono domande da proporre al committente o al proponente

        \end{itemize}
\end{itemize}

Per quanto venga fortemente incoraggiato il loro utilizzo massivo, è opportuno ribadire che lo scopo di questi strumenti è di supplemento agli incontri sincroni, non possono quindi essere utilizzato come rimpiazzo di questi ultimi. La programmazione delle attività, la suddivisione degli incarichi, la prenotazione degli incontri e tutte le altre attività di coordinamento e organizzazione del progetto rimarranno argomento delle riunioni sincrone.

\subsubsection{Riunioni Interne tra i Membri del Gruppo} \label{subsec:riunioniInterne}
Le riunioni in remoto verranno preferite a quelle in presenza in quanto la loro organizzazione è molto più semplice e si integra più facilmente con le necessità (personali e organizzative) dei componenti del gruppo. A tal proposito è stata tenuta in forte considerazione la non banale distanza geografica tra i partecipanti al progetto. \\

Verranno organizzate:
\begin{itemize}
    \item Una \textbf{riunione di aggiornamento} (tassativa) di cadenza settimanale

    \item Ulteriori \textbf{riunioni di supplemento} qualora si incontrassero difficoltà o una mole di lavoro che necessiti di un maggiore sforzo cooperativo
\end{itemize}

In entrambi i casi il giorno e l'ora degli incontri verranno decisi basandosi sul resoconto della \textbf{\underline{\hyperref[sec:pianificazioneImpegni]{disponibilità fornita}}} dai membri; al termine di ogni incontro settimanale di aggiornamento verrà programmato l'incontro successivo. Ovviamente, gli incontri settimanali di aggiornamento verranno programmati per non essere troppo vicini gli uni agli altri. Il gruppo si impegna a mantenere con rigore lo svolgimento, quantomeno, settimanale di questi incontri. Nel caso in cui si verifichino imprevisti che impediscano uno svolgimento dell'incontro tale che esso porti a un risultato utile, l'incontro verrà rinviato alla prima data libera. \\\\
Per aiutare a migliorare l'accuratezza dei verbali interni si è deciso di registrare (in formato audio) le riunioni interne. Le registrazioni verranno inviate \textbf{\underline{\hyperref[sec:canaliTelegram]{nell'apposito canale}}} del gruppo Telegram.

% DA AGGIUNGERE COME VERRANNO SVOLTE GLI INCONTRI INTERNI


\subsubsection{Comunicazioni Esterne}
Le comunicazioni con attori esterni avverranno nelle seguenti modalità:
\begin{itemize}
    \item Sincrone:
        \begin{itemize}
            \item \textbf{\underline{\nameref{strumenti:googleMeet}}}: riunioni da remoto in videochiamata 
        \end{itemize}
    \item Asincrone:
        \begin{itemize}
            \item \textbf{\underline{\hyperref[strumenti:gmail]{Email}}}: scambio di email con il committente e con il proponente tramite la email unica del gruppo
        \end{itemize}
\end{itemize}

\subsubsection{Riunioni Esterne con il Committente e il Proponente}

% DA AGGIUNGERE COME VERRANNO SVOLTE GLI INCONTRI ESTERNI


\subsection{Pianificazione degli Impegni} \label{sec:pianificazioneImpegni}
Per riuscire a facilitare la pianificazione degli impegni e tenere traccia delle disponibilità dei membri del gruppo si è deciso di usare le funzionalità offerte da \textbf{\underline{\nameref{strumenti:googleCalendar}}}. Verrà usato il calendario incorporato con la Email del gruppo, nel quale ogni membro indicherà i momenti della settimana in cui non potrà essere reperibile per eventuali incontri. Per fare ciò occorrerà che ogni componente crei un'attività nel giorno e ora in cui non sarà disponibile, inserisca come titolo dell'attività il suo nome e le dia il colore dedicato.
\\ Invece, gli impegni del gruppo come riunioni, scadenze, diari di bordo ecc. verranno contrassegnati sia aggiungendo un'attività al calendario con la rispettiva ora, sia inserendo in quel giorno un "Evento" in modo che quando si apre il calendario saltino più facilmente all'occhio i giorni in cui sono programmate attività collettive.\\\\
Le attività del calendario avranno un colore specifico in base alla loro tipologia.
Segue la legenda dei colori:

\begin{itemize}
    \item \textbf{\textcolor{blue}{Blu(o azzurro)}}: indica un momento della giornata in cui un membro del gruppo non è disponibile
    \item \textbf{\textcolor{red}{Rosso}}: indica "diari di bordo", consegne, e incontri con il committente
    \item \colorbox{gray}{\textbf{\textcolor{yellow}{Giallo}}}: indica incontri con il proponente
    \item \textbf{\textcolor{violet}{Viola}}: indica le milestone del progetto
    \item \textbf{\textcolor{green}{Verde}}: indica le riunioni interne del gruppo
\end{itemize}


\newpage
\section{Processi Primari}


\newpage
\section{Processi Secondari}

\subsection{Documentazione}

\subsubsection{Verbali Interni}
I verbali interni costituiscono un'importante risorsa, in quanto permettono di mantenere una traccia scritta degli argomenti trattati durante le riunioni interne. Contengono inoltre anche il resoconoto delle decisioni prese durante la riunione e l'assegnazione delle attività.\\
È dunque fondamentale che la stesura dei verbali sia quanto più accurata possibile. A tale scopo (\textbf{\underline{\hyperref[subsec:riunioniInterne]{come visto in precedenza}}}) è stato deciso di registrare le riunioni interne.\\\\
Per poter identificare univocamente i verbali si è deciso di usare una nomenclatura apposita; l'identificativo di ogni verbale corrisponde al nome del file. La prima lettera del nome ("V") indica che il file contiene un verbale, la seconda lettera specifica invece il tipo : "I" per i verbali interni e "E" per quelli esterni. Seguono quindi un carattere underscore ("\_") e la data della riunione oggetto del verbale, espressa con il formato anno-mese-giorno, senza caratteri di separazione.\\
La nomenclatura si conclude con il carattere underscore ("\_"), seguito dalla versione del file, espressa nel formato "vX\_Y".\\
Si è deciso di adottare un versionamento a due variabili, la prima (partendo da destra) incrementa ad ogni modifica del file, la seconda incrementa ad ogni approvazione. Ad ogni incremento della variabile X, il conteggio della variabile Y riparte da zero.\\
Esempio: "VI\_20251104\_v1\_3" indica il Verbale Interno del 4 novembre 2025, modificato tre volte dall'ultima approvazione.



\newpage
\section{Strumenti}

\subsection{Discord} \label{strumenti:discord}
\underline{\url{https://discord.com/}} \\\\
Piattaforma di messaggistica istantanea, di distribuzione digitale e di voice over IP, utilizzata per la comunicazione all'interno server Discord.
Gli utenti comunicano con chiamate vocali, videochiamate, messaggi di testo, media e file in chat private o come membri di un server Discord. Quest'ultimi sono una raccolta di canali di tipo vocale e/o testuale. 

\subsection{Git} \label{strumenti:git}
\underline{\url{https://git-scm.com/}} \\\\
Applicazione software per il controllo di versione distribuito utilizzabile da interfaccia a riga di comando. Viene utilizzato il versionamento delle repository dedicate al codice e alla documentazione del progetto.

\subsection{Github} \label{strumenti:github}
\underline{\url{https://git-scm.com/}} \\\\
È un servizio di hosting per progetti software che implementa lo strumento di controllo versione distribuito Git.  

\subsection{Gmail} \label{strumenti:gmail}
\underline{\url{https://workspace.google.com/products/gmail/}} \\\\
È un servizio gratuito di posta elettronica fornito da Google.

\subsection{Google Calendar} \label{strumenti:googleCalendar}
\underline{\url{https://workspace.google.com/products/calendar/}} \\\\
Sistema di calendari fornito da Google. Consente la gestione di un calendario condiviso, l'annotazione di attività ed eventi e la visione simultanea di tutti i calendari a cui si partecipa.


\subsection{Google Docs} \label{strumenti:googleDocs}
\underline{\url{https://workspace.google.com/products/docs/}} \\\\


\subsection{Google Meet} \label{strumenti:googleMeet}
\underline{\url{https://workspace.google.com/products/meet/}} \\\\

\subsection{Overleaf} \label{strumenti:overleaf}
\underline{\url{https://www.overleaf.com/about/features-overview}} \\\\

\subsection{Telegram} \label{strumenti:telegram}
\underline{\url{https://telegram.org/}} \\\\
Telegram è una piattaforma di messaggistica istantanea e broadcasting basato su cloud. Caratteristiche di Telegram sono la possibilità di scambiare messaggi di testo tra due utenti o tra gruppi, effettuare chiamate vocali e videochiamate cifrate point-to-point, scambiare messaggi vocali, immagini, video, file di qualsiasi tipo. Permette inoltre la suddivisione dei gruppi in canali distinti.

\subsection{WhatsApp} \label{strumenti:whatsapp}
\underline{\url{https://www.whatsapp.com/}} \\\\



\end{document}
